\section*{第 \ref{ch:g7:plate-tectonics} 章——板块构造：移动的拼图}

\begin{solution}{exo:g7:plate-tectonics:1}
\omterm{def:g7:plate-tectonics:continental-drift}{大陆漂移}是指\omterm{def:g3:continents-and-oceans:continent}{大陆}在地球表面上缓慢移动，
因此其位置随漫长的岁月而发生了变化。
魏格纳的证据包括：南美洲与非洲\omterm{def:g1:landscapes-around-us:water}{海岸}线的吻合；
如今被\omterm{def:g2:water-on-earth:ocean}{海洋}隔开的\omterm{def:g3:continents-and-oceans:continent}{大陆}上发现相同的化石；
\omterm{def:g2:water-on-earth:ocean}{海洋}两侧存在相同的\omterm{def:g3:stones-pebbles-sand:rock}{岩石}和古老的山链；古\omterm{def:g5:weather-and-climate:climate}{气候}遗迹
位置不合常理（\omterm{def:g3:continents-and-oceans:poles}{南极}洲有煤层，炎热干燥的非洲有冰川\omterm{def:g5:water-shapes-the-land:deposition}{沉积}）。
\end{solution}

\begin{solution}{exo:g7:plate-tectonics:2}
离散型：板块相互远离。汇聚型：板块彼此靠近。
转换型：板块彼此错动、并排滑过。
\end{solution}

\begin{solution}{exo:g7:plate-tectonics:3}
例如，环太平洋地区（日本、印度尼西亚、安第斯\omterm{def:g4:world-physical-map:range}{山脉}）以及
从地中海延伸至喜马拉雅山的\omterm{def:g4:volcanoes-earthquakes-first-look:earthquake}{地震}带。非洲腹地和澳大利亚
几乎没有大\omterm{def:g4:volcanoes-earthquakes-first-look:earthquake}{地震}（洛基\omterm{def:g4:world-physical-map:range}{山脉}以东的北美洲，
或巴西，也同样如此）。
\end{solution}

\begin{solution}{exo:g7:plate-tectonics:4}
\omterm{def:g2:water-on-earth:salt-fresh}{淡水}动物无法穿越数千千米
的\omterm{def:g2:water-on-earth:salt-fresh}{咸水}\omterm{def:g2:water-on-earth:ocean}{海洋}。两侧都发现了它，说明两地
曾在 \emph{Mesosaurus} 生活的时代同属一块\omterm{def:g3:continents-and-oceans:continent}{大陆}、以陆地相连。
\end{solution}

\begin{solution}{exo:g7:plate-tectonics:5}
煤由温暖、潮湿的沼泽森林中的植物形成：当时\omterm{def:g5:weather-and-climate:climate}{气候}温暖湿润。
如今这类森林无法在\omterm{def:g3:continents-and-oceans:poles}{南极}附近生长，因此\omterm{def:g3:continents-and-oceans:poles}{南极}洲
过去一定离\omterm{def:g3:continents-and-oceans:equator}{赤道}近得多：它已经发生了移动。
\end{solution}

\begin{solution}{exo:g7:plate-tectonics:6}
大西洋中脊：离散型。安第斯\omterm{def:g4:world-physical-map:range}{山脉}：汇聚型（大洋板块
俯冲至\omterm{def:g3:continents-and-oceans:continent}{大陆}之下）。圣安德烈亚斯断层：转换型。喜马拉雅山：汇聚型（两个
\omterm{def:g3:continents-and-oceans:continent}{大陆}相撞）。东非裂谷：离散型（\omterm{def:g3:continents-and-oceans:continent}{大陆}正在裂开）。
\end{solution}

\begin{solution}{exo:g7:plate-tectonics:7}
$2.5 \times 10 = 25$ 千米每百万年；经过 6000 万年：
$25 \times 60 = \qty{1500}{km}$。
\end{solution}

\begin{solution}{exo:g7:plate-tectonics:8}
$\qty{7}{m} = \qty{700}{cm}$，$700 \div 9 \approx 78$：约
\qty{78}{cm} 每年，约为平均值 \qty{2.5}{cm} 每年的 30 倍。
板块在地表并非平稳分离：应变会累积数年，
随后以短促的爆发释放，裂缝一次张开数米。
长期来看，平静年份与爆发期平均下来，
每年约为几厘米。
\end{solution}

\begin{solution}{exo:g7:plate-tectonics:9}
魏格纳无法说明是什么力推动\omterm{def:g3:continents-and-oceans:continent}{大陆}，他所设想的运动
（\omterm{def:g3:continents-and-oceans:continent}{大陆}在洋底中犁行）是不可能发生的。
在\omterm{def:g7:plate-tectonics:plate-tectonics}{板块构造}理论中，\omterm{def:g3:continents-and-oceans:continent}{大陆}并不穿行于任何东西之中：它们
与周围的洋底一起被携带，携带它们的是作为地球外壳整块部分
而运动的板块。
\end{solution}

\begin{solution}{exo:g7:plate-tectonics:10}
洋脊处最高，海沟处最深。新板块在洋脊处形成，
热\omterm{def:g3:stones-pebbles-sand:rock}{岩石}从\omterm{def:g6:inside-the-earth:layers}{地幔}上涌至此；大洋板块在海沟处
重新没入\omterm{def:g6:inside-the-earth:layers}{地幔}，俯冲至\omterm{def:g3:continents-and-oceans:continent}{大陆}之下。
\end{solution}

\begin{solution}{exo:g7:plate-tectonics:11}
$14.2 \div 2.2 \approx 6.45$：约快 6 倍。
以千米每百万年计，速率分别为 142 和 22。行进 \qty{1000}{km}
所需时间：东太平洋海隆 $1000 \div 142 \approx 7$，约 700 万年；
北大西洋中脊 $1000 \div 22 \approx 45$，约
4500 万年。
\end{solution}

\begin{solution}{exo:g7:plate-tectonics:12}
(a) $6472 - 4578 = 1894$ 次\omterm{def:g4:volcanoes-earthquakes-first-look:earthquake}{地震}。(b) $4578 \div 7031 \approx 0.651$，
即 65\,\%。(c) 图上的线标示海沟，大洋
板块开始下沉。板块并非竖直向下：它
倾斜地俯冲到另一板块之下，越往下越深。
\omterm{def:g4:volcanoes-earthquakes-first-look:earthquake}{地震}沿下沉的板块处处发生，因此许多\omterm{def:g4:volcanoes-earthquakes-first-look:earthquake}{地震}位于
上覆板块之下，与海沟处于同一侧，距该线数十
至数百千米。\omterm{def:g7:plate-tectonics:plate-boundary}{板块边界}是一条带，而不是一条细线。
\end{solution}

\begin{solution}{pb:g7:plate-tectonics:1}
\textbf{1.} \emph{Mesosaurus}、\emph{Cynognathus} 和 \emph{Glossopteris}
（\emph{Lystrosaurus} 的化石带连接非洲、印度和\omterm{def:g3:continents-and-oceans:poles}{南极}洲）。

\textbf{2.} 它生活在\omterm{def:g2:water-on-earth:salt-fresh}{淡水}中，无法穿越
数千千米宽的\omterm{def:g2:water-on-earth:salt-fresh}{咸水}\omterm{def:g2:water-on-earth:ocean}{海洋}。

\textbf{3.} 巴西东北部的凸出部分恰好嵌入
非洲西\omterm{def:g1:landscapes-around-us:water}{海岸}的几内亚湾凹陷处。

\textbf{4.} \omterm{def:g7:plate-tectonics:plate-boundary}{离散型边界}，即大西洋中脊：南美板块
向西移动，非洲板块向东移动，两者相互远离。

\textbf{5.} 两侧\omterm{def:g1:landscapes-around-us:water}{海岸}分处\omterm{def:g6:latitude-longitude:longitude}{本初子午线}两侧：
$43.2 + 14.4 = \ang{57.6}$。

\textbf{6.} $57.6 \times 102 = 5875.2$：约 \qty{5880}{km}。

\textbf{7.} 从\omterm{def:g3:continents-and-oceans:equator}{赤道}向极点移动时，\omterm{def:g6:latitude-longitude:longitude}{经线}
越来越靠近。南纬 \ang{23} \omterm{def:g6:latitude-longitude:latitude}{纬线}圈比\omterm{def:g3:continents-and-oceans:equator}{赤道}小，
因此其 360 度中的每一度都更短。

\textbf{8.} 约为宽度的一半：$5880 \div 2 = \qty{2940}{km}$。

\textbf{9.} $3.4 \times 10 = 34$ 千米每百万年。

\textbf{10.} $5880 \div 34 \approx 173$：约 1.7 亿年。

\textbf{11.} 是的：1.7 亿年介于 1 亿年与 2 亿年之间，
并且晚于\omterm{def:g7:plate-tectonics:supercontinent}{泛大陆}开始解体的时间（2.25 亿至
2 亿年前）。在\omterm{def:g2:drawing-to-map:map}{地图}上做简单计算即可得到
正确的数量级。

\textbf{12.} 例如：速率并非始终相同；\omterm{def:g3:continents-and-oceans:continent}{大陆}在
分裂前已经拉伸变薄，因此今天的\omterm{def:g1:landscapes-around-us:water}{海岸}线
并不恰好是它们当初断裂的那条线。
\end{solution}
